\documentclass[conference]{IEEEtran}
\IEEEoverridecommandlockouts

\usepackage{cite}
\usepackage{amsmath,amssymb,amsfonts}
\usepackage{graphicx}
\usepackage{textcomp}
\usepackage{xcolor}
\usepackage{booktabs}
\usepackage{multirow}
\usepackage{makecell}
\usepackage{url}

\begin{document}

\title{Face Authentication Using Real vs. Synthetic\\Training Data:\\
A Comparative Study Based on NIST FRTE\\Evaluation Metrics\\[0.4em]
{\large \textit{Zwischenbericht / Interim Research Report}}}

\author{
  \IEEEauthorblockN{Nour Ebid}
  \IEEEauthorblockA{
    Seminar zu aktuellen Entwicklungen\\
    M.Sc.\ Angewandte Informatik, Fachbereich 4\\
    Hochschule f{\"u}r Technik und Wirtschaft Berlin\\
    Berlin, Deutschland\\
    Nour.Ebid@Student.HTW-Berlin.de
  }
}

\maketitle

% ----------------------------------------------------------------
\begin{abstract}
Face recognition is a widely deployed biometric authentication
technology, yet its reliance on real facial photographs raises
significant privacy, consent, and data availability concerns.
Under the General Data Protection Regulation (GDPR), biometric
data are classified as sensitive personal data, making the
storage of real facial images a legal and logistical challenge
for organisations deploying access control systems. This
interim report introduces an ongoing study that investigates
whether AI-generated synthetic facial images can replace real
photographs entirely during the registration phase of a face
authentication system, without a significant loss in security
performance. Two experiments are being designed under identical
evaluation conditions aligned with the NIST Face Recognition
Technology Evaluation (FRTE) framework: Experiment~1
establishes a real-data baseline, and Experiment~2 registers
the same users using synthetically generated images. This
report presents the motivation, background, research
hypothesis, and experimental variables. Methodology and
results will follow in the final paper.
\end{abstract}

\begin{IEEEkeywords}
face authentication, synthetic data, FaceNet, NIST FRTE,
false accept rate, biometric verification, privacy-preserving
biometrics, identity-preserving augmentation
\end{IEEEkeywords}

% ================================================================
\section{Introduction}

Biometric authentication systems based on facial recognition
have become a cornerstone of modern digital security, deployed
in applications ranging from smartphone unlock to airport
border control. Unlike password-based systems, face
authentication offers passive, non-intrusive verification
without requiring the user to remember credentials. These
systems fundamentally rely on databases of registered facial
images to calibrate recognition models for each authorised
user.

The collection and storage of real facial photographs,
however, introduces a set of practical and regulatory
challenges that are increasingly difficult to ignore. Under
the European Union's General Data Protection Regulation
(GDPR)~\cite{gdpr}, biometric data---including facial
images---are classified as a special category of sensitive
personal data under Article~9, requiring explicit informed
consent, purpose limitation, and strict storage controls.
Organisations deploying face authentication must therefore
store, secure, and manage real biometric data, exposing them
to data breach liability and ongoing compliance obligations.
Beyond regulatory concerns, reproducing consistent imaging
conditions across enrollment sessions is logistically
difficult: ambient lighting, camera type, and subject pose
vary freely in naturalistic settings, introducing systematic
variability that degrades authentication reliability.

AI-generated synthetic facial images offer a compelling
alternative. Advances in Generative Adversarial
Networks~\cite{goodfellow2014} and diffusion-based synthesis
models~\cite{rombach2022} enable the creation of
photorealistic human faces without storing sensitive biometric
originals. If a face authentication system registered
exclusively on synthetic representations of users performs
comparably to one registered on real photographs, then
organisations could enroll users using synthetic data
only---eliminating the need to retain real biometric images
altogether. This would represent a significant step toward
privacy-preserving biometric authentication.

The central research question addressed in this study is:

\begin{quote}
  \textit{Can a face authentication system enrolled exclusively
  on AI-generated synthetic facial images achieve functionally
  acceptable security performance comparable to a system
  enrolled on real photographs, making synthetic data a
  viable privacy-preserving alternative?}
\end{quote}

To evaluate this, we implement a face authentication pipeline
based on the FaceNet deep neural network~\cite{schroff2015},
register four authorised users under both real and synthetic
conditions, and evaluate both systems using False Accept
Rate~(FAR), False Reject Rate~(FRR), and accuracy on an
identical held-out test set, following the NIST FRTE
evaluation framework~\cite{nist2024}.

The remainder of this report is structured as follows.
Section~II reviews relevant background and related work.
Section~III formalises the research hypothesis and
experimental variables.

% ================================================================
\section{Background and Related Work}

\subsection{Face Verification vs.\ Identification}

Face recognition covers two distinct operational tasks.
\textit{Identification} (closed-set or open-set) answers ``who
is this person?'' by searching a database of known identities
for the best match. \textit{Verification} (authentication)
answers a narrower question: ``is this person who they claim
to be?'' It performs a one-to-one comparison between a query
embedding and a stored reference for the claimed identity.
This study addresses the verification task exclusively, which
is the relevant subproblem in access control systems where the
claimed identity is known a priori.

\subsection{FaceNet Embedding-Based Verification}

FaceNet~\cite{schroff2015} is a deep convolutional neural
network trained using a triplet loss to produce a
128-dimensional Euclidean embedding space in which intra-person
distances are minimised and inter-person distances are
maximised. For an anchor, positive, and negative sample
$(x^a, x^p, x^n)$, the triplet objective is:
\begin{equation}
  \|f(x^a) - f(x^p)\|_2^2 + \alpha
    < \|f(x^a) - f(x^n)\|_2^2
\end{equation}
where $f(\cdot)$ denotes the learned embedding function and
$\alpha$ is the margin. Verification is performed by computing
the L2 distance between a query embedding $\mathbf{e}_q$ and
a stored reference embedding $\bar{\mathbf{e}}_u$:
\begin{equation}
  d = \|\mathbf{e}_q - \bar{\mathbf{e}}_u\|_2
\end{equation}
\begin{equation}
  \text{Decision} =
    \begin{cases}
      \text{VERIFIED}  & d < \theta \\
      \text{IMPOSTOR}  & d \geq \theta
    \end{cases}
\end{equation}
FaceNet is accessed via the DeepFace library~\cite{serengil2020}
with pre-trained weights, without retraining on experimental
subjects. This makes it suitable as a fixed, model-controlled
comparison baseline.

\subsection{NIST FRTE Evaluation Metrics}

The NIST Face Recognition Technology Evaluation
(FRTE)~\cite{nist2024} establishes the standard evaluation
protocol for biometric face recognition systems. The primary
metrics used in this study are:
\begin{itemize}
  \item \textbf{FAR (False Accept Rate):} the fraction of
        impostor attempts incorrectly accepted as genuine.
        A high FAR represents a direct security vulnerability.
  \item \textbf{FRR (False Reject Rate):} the fraction of
        genuine user attempts incorrectly rejected. A high FRR
        degrades usability by inconveniencing legitimate users.
  \item \textbf{Accuracy:} the overall fraction of correct
        decisions across both genuine and impostor queries.
\end{itemize}
A confusion matrix (TP, FP, FN, TN) is recorded for each
experiment to support transparent, reproducible evaluation.

\subsection{Privacy Challenges in Biometric Enrollment}

The GDPR~\cite{gdpr} classifies biometric data as a special
category of personal data under Article~9, imposing strict
requirements on collection, storage, and processing. In the
context of face authentication, this means every registered
user's real facial photographs must be handled under explicit
consent and data protection controls. A system that achieves
comparable security without storing any real photographs would
satisfy these requirements by design---reducing both
compliance burden and breach exposure.

\subsection{Synthetic Data in Face Recognition Research}

Prior work has explored synthetic data as a substitute for
real data in face recognition training. Wood et
al.~\cite{wood2021} demonstrated that a face analysis model
trained exclusively on 3D-rendered synthetic faces---with no
real training images---achieves competitive performance on
facial landmark detection benchmarks. Bae et al.~\cite{bae2023}
introduced DigiFace-1M, one million digitally rendered
synthetic face images, used to train recognition models
achieving 95.4\% LFW accuracy. Both works, however, focus on
the model \textit{training} (identification) setting. This
study targets the \textit{enrollment} (verification) setting:
whether synthetic images of a known user can serve as the
registered reference, directly replacing real photographs in a
deployed access control system. This distinction is more
practically significant and remains underexplored.

% ================================================================
\section{Research Hypothesis and Variables}

\subsection{Hypothesis}

This study adopts a privacy-preserving framing. The core
question is not whether real data \textit{outperforms}
synthetic data---it is expected to, given the direct
correspondence between enrollment and test conditions---but
whether the synthetic system achieves \textit{functionally
acceptable} performance that would justify replacing real
data in practice.

The research hypothesis is therefore:

\begin{quote}
  \textbf{H$_1$:} \textit{A face authentication system
  enrolled exclusively on AI-generated synthetic facial images
  of registered users can achieve functionally acceptable
  security performance---specifically FAR $\leq 15\%$ and
  FRR $\leq 5\%$---making it a viable privacy-preserving
  alternative to a system enrolled on real photographs of
  the same individuals.}
\end{quote}

\begin{quote}
  \textbf{H$_0$:} \textit{The synthetic-enrolled system
  fails to meet the defined performance thresholds
  (FAR $> 15\%$ or FRR $> 5\%$), meaning synthetic data
  is not a reliable substitute for real biometric enrollment
  data under these conditions.}
\end{quote}

The 15\% FAR and 5\% FRR thresholds are operationally
motivated: a FAR above 15\% represents an unacceptable
security risk for a real access control system, while an
FRR above 5\% would cause unacceptable inconvenience to
legitimate users. These bounds are consistent with tolerance
margins observed in NIST FRTE evaluations across
demographically similar systems~\cite{nist2024}.

If H$_1$ is confirmed, it provides empirical evidence that
organisations can eliminate real biometric data storage from
their enrollment pipeline without meaningfully compromising
security---a directly actionable privacy-preserving finding.

\subsection{Variables}

Table~\ref{tab:variables} summarises all experimental
variables. Strict control of the \textit{Controlled}
variables is essential: any observed difference in FAR or FRR
between the two experiments must be attributable solely to
the nature of the enrollment data (real vs.\ synthetic), not
to differences in the recognition model, threshold, test set
composition, or evaluation protocol.

\begin{table}[htbp]
  \caption{Experimental Variables}
  \label{tab:variables}
  \centering
  \begin{tabular}{lll}
    \toprule
    \textbf{Type} & \textbf{Variable} & \textbf{Value / Range} \\
    \midrule
    Independent & Enrollment data type
      & Real photos vs.\ synthetic images \\
    \midrule
    \multirow{3}{*}{Dependent}
      & FAR      & Measured (\%) \\
      & FRR      & Measured (\%) \\
      & Accuracy & Measured (\%) \\
    \midrule
    \multirow{5}{*}{Controlled}
      & Model              & FaceNet (identical) \\
      & Threshold $\theta$ & Fixed (identical) \\
      & Enrollment images  & 10 per person \\
      & Test images        & 5 held-out per person \\
      & Impostors          & 20 identical AI strangers \\
    \midrule
    \multirow{4}{*}{\makecell{Controlled\\(synthesis)}}
      & Pose      & Canonical alignment \\
      & Lighting  & Gamma correction \\
      & Skin tone & CLAHE in LAB space \\
      & Camera    & Blur + noise \\
    \bottomrule
  \end{tabular}
\end{table}

% ================================================================
\begin{thebibliography}{9}

\bibitem{schroff2015}
F.~Schroff, D.~Kalenichenko, and J.~Philbin,
``FaceNet: A unified embedding for face recognition and
clustering,''
\textit{Proc. IEEE CVPR}, pp.~815--823, 2015.

\bibitem{nist2024}
National Institute of Standards and Technology,
``Face Recognition Technology Evaluation (FRTE),''
\url{https://www.nist.gov/programs-projects/face-recognition-technology-evaluation-frte},
2024.

\bibitem{serengil2020}
S.~I.~Serengil and A.~Ozpinar,
``LightFace: A hybrid deep face recognition framework,''
\textit{Proc. IEEE ASYU}, 2020.

\bibitem{wood2021}
E.~Wood et al.,
``Fake it till you make it: Face analysis in the wild using
synthetic data alone,''
\textit{Proc. IEEE ICCV}, 2021.

\bibitem{bae2023}
G.~Bae et al.,
``DigiFace-1M: 1 million digital face images for face
recognition,''
\textit{Proc. IEEE WACV}, 2023.

\bibitem{gdpr}
European Parliament and Council,
``Regulation (EU) 2016/679 --- General Data Protection
Regulation (GDPR),''
\textit{Official Journal of the European Union},
L~119, pp.~1--88, 2016.

\bibitem{goodfellow2014}
I.~Goodfellow et al.,
``Generative adversarial nets,''
\textit{Adv. Neural Inf. Process. Syst.}, vol.~27, 2014.

\bibitem{rombach2022}
R.~Rombach, A.~Blattmann, D.~Lorenz, P.~Esser, and B.~Ommer,
``High-resolution image synthesis with latent diffusion
models,''
\textit{Proc. IEEE CVPR}, pp.~10684--10695, 2022.

\end{thebibliography}

\end{document}
